% Hanzo AI Cloud: Confidential Inference, Agent Hosting, and Model Serving
\documentclass[11pt]{article}
\usepackage{shared/hanzocover}
\usepackage{amsmath, amssymb, amsthm}
\usepackage{mathtools}
\usepackage{bm}
\usepackage{graphicx}
\usepackage{booktabs}
\usepackage{hyperref}
\usepackage{enumitem}
\usepackage{algorithm}
\usepackage{algpseudocode}
\usepackage{xcolor}
\usepackage{float}
\usepackage{multirow}
\usepackage{listings}
\input{shared/lstlang}
\usepackage{tikz}
\usetikzlibrary{shapes,arrows,positioning,fit}
\usepackage{pgfplots}
\pgfplotsset{compat=1.18}

\hypersetup{colorlinks=true,linkcolor=black,citecolor=blue,urlcolor=blue}

\definecolor{codegreen}{rgb}{0,0.6,0}
\definecolor{codegray}{rgb}{0.5,0.5,0.5}
\definecolor{codepurple}{rgb}{0.58,0,0.82}
\definecolor{backcolour}{rgb}{0.95,0.95,0.92}

\lstdefinestyle{mystyle}{
    backgroundcolor=\color{backcolour},
    commentstyle=\color{codegreen},
    keywordstyle=\color{codepurple},
    numberstyle=\tiny\color{codegray},
    stringstyle=\color{codegreen},
    basicstyle=\ttfamily\footnotesize,
    breakatwhitespace=false,
    breaklines=true,
    captionpos=b,
    keepspaces=true,
    numbers=left,
    numbersep=5pt,
    showspaces=false,
    showstringspaces=false,
    showtabs=false,
    tabsize=2
}
\lstset{style=mystyle}

\newtheorem{theorem}{Theorem}
\newtheorem{definition}{Definition}

\title{HANZO AI CLOUD\\[0.3em]
{\large Confidential Inference, Agent Hosting, and\\Model Serving on Sovereign Infrastructure}}
\author{
    Zach Kelling \\
    \textit{Hanzo Industries Inc (Techstars '17)} \\
    \texttt{research@hanzo.ai}
}
\date{March 2026}

\begin{document}
\hanzocoverpage

% ============================================================================
% ABSTRACT
% ============================================================================
\begin{abstract}
We present Hanzo AI Cloud, an infrastructure layer for AI workloads where the
provider never sees the user's data. Current AI cloud platforms (OpenAI,
Anthropic, Replicate, Together) require users to transmit plaintext prompts,
receive plaintext completions, and trust the provider not to log, train on, or
leak their inputs. This trust model is incompatible with regulated industries
(healthcare, finance, defense) and with any user who values privacy as a
technical guarantee rather than a policy promise.

Hanzo AI Cloud eliminates this trust requirement through three mechanisms:
(i)~confidential inference via CKKS fully homomorphic encryption, where inputs
are encrypted client-side, evaluated homomorphically on the provider's GPUs,
and decrypted only by the client; (ii)~agent hosting with capsule isolation,
where each agent receives a private vault, a decentralized identifier, a
capability budget, and encrypted memory that no other agent or the platform
can access; and (iii)~model serving with verifiable attestation, where every
inference result is accompanied by a cryptographic receipt anchored to the Lux
A-Chain, proving which model produced it, on what hardware, at what time.

We describe the system architecture, the FHE inference pipeline, the agent
capsule model, the model registry, Hanzo Candle (the Rust ML framework
powering GPU inference), the Jin multimodal architecture, edge inference for
mobile and IoT devices, and the provider marketplace with stake-weighted
selection and SLA enforcement.
\end{abstract}

% ============================================================================
\section{Introduction}
\label{sec:intro}
% ============================================================================

Every major AI cloud platform operates on the same trust model: the user sends
plaintext data to the provider, the provider runs a model on it, the provider
returns a plaintext result. The provider sees everything. The provider's
privacy policy is the only barrier between the user's data and the provider's
training pipeline, the provider's employees, the provider's government
subpoenas, and the provider's security breaches.

This is not a theoretical concern. OpenAI's data retention policy permits
logging all API inputs for 30 days~\cite{openaiapi}. Anthropic retains inputs
for safety evaluation. Replicate processes inputs on shared GPU instances with
no hardware isolation between tenants. The user has no cryptographic guarantee
that their prompt---containing patient records, financial models, legal
strategy, classified intelligence---is not stored, read, or reused.

Hanzo AI Cloud solves this problem at the architectural level, not the policy
level. The system is designed so that the provider \emph{cannot} see the
user's data, not that the provider \emph{promises} not to.

\subsection{Contributions}

\begin{enumerate}[leftmargin=1.4em]
    \item Confidential inference via CKKS fully homomorphic encryption with
    threshold decryption (\S\ref{sec:fhe}).
    \item Agent hosting with capsule isolation: vault, DID, capabilities,
    budget, and private memory (\S\ref{sec:agents}).
    \item Model serving with verifiable inference receipts anchored to the
    Lux A-Chain (\S\ref{sec:serving}).
    \item Hanzo Candle: a Rust ML framework with GPU acceleration integrated
    into Base Functions (\S\ref{sec:candle}).
    \item Jin architecture: unified multimodal inference across vision,
    language, and audio (\S\ref{sec:jin}).
    \item Edge inference: local models on mobile/IoT with CRDT result sync
    (\S\ref{sec:edge}).
    \item Provider marketplace with stake-weighted selection and on-chain SLA
    enforcement (\S\ref{sec:marketplace}).
\end{enumerate}

% ============================================================================
\section{The AI Cloud Problem}
\label{sec:problem}
% ============================================================================

The fundamental issue is that inference requires the model to operate on
plaintext. A transformer model computes attention over token embeddings; if
the embeddings are encrypted, standard matrix multiplication does not apply.
This creates an apparent dilemma: either the provider sees the data (plaintext
inference) or the provider cannot run the model (encrypted data).

Three approaches exist to break this dilemma:

\begin{enumerate}[leftmargin=1.4em]
    \item \textbf{Trusted Execution Environments (TEEs)}: Run inference inside
    an enclave (Intel SGX, AMD SEV, ARM TrustZone). The hardware prevents the
    host from reading enclave memory. Limitation: TEEs have small memory
    (256~MB for SGX), side-channel vulnerabilities, and require trusting the
    hardware vendor's attestation.
    \item \textbf{Secure Multi-Party Computation (MPC)}: Split the computation
    across multiple non-colluding parties. No single party sees the full
    input. Limitation: communication overhead scales with model size,
    making MPC impractical for billion-parameter models.
    \item \textbf{Fully Homomorphic Encryption (FHE)}: Encrypt the input and
    evaluate the model on ciphertexts. The provider performs computation
    without decryption. Limitation: computational overhead (10--1000$\times$
    slowdown depending on scheme and circuit depth).
\end{enumerate}

Hanzo AI Cloud uses FHE as the primary confidentiality mechanism, with TEEs as
an optional accelerator for latency-sensitive workloads. MPC is used only for
threshold decryption of FHE outputs, where the communication overhead is
minimal (one round, fixed-size messages).

% ============================================================================
\section{Confidential Inference}
\label{sec:fhe}
% ============================================================================

Hanzo implements confidential inference using the CKKS (Cheon--Kim--Kim--Song)
FHE scheme~\cite{ckks2017}, which supports approximate fixed-point arithmetic
on encrypted vectors---the operation profile that neural network inference
requires.

\subsection{CKKS Overview}

CKKS operates on the ring $R_Q = \mathbb{Z}_Q[X]/(X^N + 1)$ where $N$ is a
power of 2 (typically $N = 2^{16}$) and $Q$ is a product of coprime moduli.
A plaintext vector $\bm{z} \in \mathbb{C}^{N/2}$ is encoded as a ring
element $m \in R_Q$ via the canonical embedding, then encrypted as:

\begin{equation}
\text{ct} = (\bm{a} \cdot s + e + m, \; \bm{a})
\end{equation}

where $s$ is the secret key, $\bm{a}$ is uniform random, and $e$ is a small
error. Homomorphic addition and multiplication on ciphertexts correspond to
addition and multiplication on plaintexts, with the error growing after each
operation.

\subsection{Inference Pipeline}

The confidential inference pipeline operates in four phases:

\begin{enumerate}[leftmargin=1.4em]
    \item \textbf{Client-side encryption}: The client encodes the input (token
    embeddings for text, pixel values for images) into CKKS plaintext vectors
    and encrypts them with their public key. The ciphertexts are transmitted
    to the Hanzo Cloud.
    \item \textbf{Homomorphic evaluation}: The cloud evaluates the model on
    the ciphertexts. Linear layers (matrix multiplication, convolution) map
    directly to CKKS operations. Non-linear activations (GELU, softmax) are
    approximated by low-degree polynomials fitted via Remez
    algorithm~\cite{hanzocandle, hanzofhe}.
    \item \textbf{Encrypted output}: The cloud returns the encrypted result.
    It cannot decrypt it---only the client (or threshold decryption
    participants) hold the secret key.
    \item \textbf{Client-side decryption}: The client decrypts the result
    locally. For threshold configurations, $k$-of-$n$ key holders cooperate
    to produce the decryption.
\end{enumerate}

\subsection{Polynomial Activation Approximation}

The GELU activation $\text{GELU}(x) = x \cdot \Phi(x)$ where $\Phi$ is the
Gaussian CDF is approximated by a degree-7 polynomial on $[-8, 8]$:

\begin{equation}
\widetilde{\text{GELU}}(x) = \sum_{k=0}^{7} a_k x^k
\end{equation}

with coefficients fitted to minimize $L_\infty$ error. The maximum
approximation error is $|\text{GELU}(x) - \widetilde{\text{GELU}}(x)| <
2 \times 10^{-4}$ on the domain, which is within the noise floor of CKKS
arithmetic. Softmax is decomposed into exponentiation (degree-6 polynomial
approximation) and reciprocal (Goldschmidt iteration on ciphertexts).

\subsection{Threshold Decryption}

For multi-party scenarios, the secret key $s$ is split into $n$ shares
$s_1, \ldots, s_n$ via Shamir's secret sharing. Decryption requires $k$
shares:

\begin{equation}
m = \text{ct}_0 + \text{ct}_1 \cdot s = \text{ct}_0 +
\text{ct}_1 \cdot \sum_{i \in S} \lambda_i s_i
\end{equation}

where $S \subset [n]$ with $|S| = k$ and $\lambda_i$ are Lagrange
interpolation coefficients. Each participant computes a partial decryption
$d_i = \text{ct}_1 \cdot \lambda_i s_i + e_i$ (with a smudging noise $e_i$
to prevent key leakage), and the combiner sums the partial decryptions. No
single party learns the plaintext unless they hold $k$ shares.

% ============================================================================
\section{Agent Hosting}
\label{sec:agents}
% ============================================================================

Hanzo AI Cloud hosts autonomous agents---persistent AI processes with memory,
tools, and goals. Each agent runs inside a \textbf{capsule}: an isolated
execution environment with its own data, identity, and resource budget.

\subsection{Capsule Architecture}

\begin{definition}[Agent Capsule]
A capsule is a tuple $(\text{vault}, \text{did}, \text{caps}, \text{budget},
\text{mem})$ where:
\begin{itemize}[leftmargin=1.1em]
    \item $\text{vault}$: A private Base instance (encrypted SQLite) containing
    the agent's persistent state.
    \item $\text{did}$: A decentralized identifier on the Lux I-Chain,
    uniquely identifying the agent.
    \item $\text{caps}$: A capability set defining what the agent can access
    (collections, APIs, other agents, network).
    \item $\text{budget}$: A HANZO token budget limiting the agent's compute
    and API spending.
    \item $\text{mem}$: Encrypted working memory (conversation history,
    scratchpad, retrieved context) stored in the vault.
\end{itemize}
\end{definition}

\subsection{Isolation Guarantees}

Capsule isolation is enforced at three levels:

\begin{enumerate}[leftmargin=1.4em]
    \item \textbf{Storage isolation}: Each capsule's vault is a separate
    encrypted SQLite database with a unique encryption key. The platform
    cannot read the vault contents without the agent's key.
    \item \textbf{Compute isolation}: Each capsule runs in a separate Goja VM
    instance (for tool execution) and a separate inference context (for model
    calls). No shared state between capsules.
    \item \textbf{Network isolation}: Capsule network access is restricted to
    the capabilities declared in $\text{caps}$. An agent without the
    \texttt{network} capability cannot make outbound requests.
\end{enumerate}

\subsection{Agent Memory}

Agent memory is structured as three tiers:

\begin{table}[H]
\centering
\small
\begin{tabular}{llp{5cm}}
\toprule
\textbf{Tier} & \textbf{Storage} & \textbf{Behavior} \\
\midrule
Working & In-process & Current conversation context. Lost on capsule
restart. \\
Episodic & Vault (encrypted) & Past interactions, indexed by time and topic.
Persists across restarts. Searchable via vector similarity. \\
Semantic & Vault (encrypted) & Distilled knowledge extracted from episodic
memory. Compact, structured, versioned. \\
\bottomrule
\end{tabular}
\caption{Agent memory tiers. All persistent memory is encrypted in the
capsule vault.}
\label{tab:memory}
\end{table}

Episodic-to-semantic distillation runs as a scheduled Base Function inside the
capsule. The agent periodically reviews its episodic memory and extracts
durable facts, preferences, and learned patterns into the semantic tier. This
is analogous to sleep consolidation in biological memory systems.

\subsection{Budget Enforcement}

Each capsule has a HANZO token budget that limits:

\begin{itemize}[leftmargin=1.1em]
    \item \textbf{Inference tokens}: Maximum tokens generated per period.
    \item \textbf{API calls}: Maximum outbound API calls per period.
    \item \textbf{Storage}: Maximum vault size in bytes.
    \item \textbf{Compute}: Maximum CPU-seconds for function execution.
\end{itemize}

Budget is enforced at the platform level. When a capsule exhausts its budget,
it is suspended (not terminated). The capsule's vault and memory persist.
The owner can replenish the budget to resume the agent.

% ============================================================================
\section{Model Serving}
\label{sec:serving}
% ============================================================================

Hanzo AI Cloud serves models from a registry anchored to the Lux A-Chain
(Attestation Chain). Every model version, every inference result, and every
provider claim is recorded on-chain with cryptographic receipts.

\subsection{Model Registry}

The model registry stores model metadata on the A-Chain:

\begin{table}[H]
\centering
\small
\begin{tabular}{lp{6cm}}
\toprule
\textbf{Field} & \textbf{Description} \\
\midrule
\texttt{model\_id} & Content-addressed hash of model weights \\
\texttt{format} & ONNX, GGUF, SafeTensors, Candle \\
\texttt{architecture} & Transformer, CNN, Diffusion, MoE \\
\texttt{parameters} & Total parameter count \\
\texttt{quantization} & FP16, INT8, INT4, or none \\
\texttt{benchmarks} & Standardized evaluation scores \\
\texttt{publisher} & I-Chain DID of the publishing entity \\
\texttt{license} & Machine-readable license identifier \\
\bottomrule
\end{tabular}
\caption{Model registry fields. The \texttt{model\_id} is a SHA3-256 hash of
the model weights, ensuring immutable identification.}
\label{tab:registry}
\end{table}

Model weights are stored off-chain (on Hanzo Storage, S3-compatible) with the
hash anchored on-chain. This ensures that a model ID always refers to the same
weights---the provider cannot silently swap a model.

\subsection{Verifiable Inference Receipts}

Every inference call produces a receipt:

\begin{equation}
\text{receipt} = \text{sign}_{k_{\text{provider}}}\bigl(
\text{model\_id} \| \text{input\_hash} \| \text{output\_hash} \|
\text{hw\_id} \| \text{timestamp}
\bigr)
\end{equation}

The receipt is submitted to the A-Chain. A client can verify: (i) which model
produced the output, (ii) on what hardware, (iii) at what time, and (iv) that
the provider's stake is intact (not slashed for prior misbehavior). Receipts
enable auditability for regulated industries where ``an AI said so'' is not
sufficient---the specific model version and provider must be traceable.

\subsection{GPU Scheduling}

The Hanzo scheduler assigns inference requests to providers based on:

\begin{itemize}[leftmargin=1.1em]
    \item \textbf{Model availability}: Does the provider have the requested
    model loaded in GPU memory?
    \item \textbf{Latency estimate}: Predicted latency based on provider
    location, current load, and historical performance.
    \item \textbf{Stake weight}: Providers with higher HANZO stake receive
    priority. Stake is a credible commitment to quality.
    \item \textbf{Cost}: The provider's advertised price per token/image/second.
\end{itemize}

The scheduler uses a modified weighted round-robin that favors loaded models
(avoiding cold-load latency) while maintaining fairness across providers.

% ============================================================================
\section{Hanzo Candle}
\label{sec:candle}
% ============================================================================

Hanzo Candle is a Rust ML framework purpose-built for inference
workloads~\cite{hanzocandle}. It provides the GPU acceleration layer for
Hanzo AI Cloud.

\subsection{Design Principles}

\begin{itemize}[leftmargin=1.1em]
    \item \textbf{Inference-first}: Candle optimizes for inference, not
    training. The tensor graph is static (no autograd by default), enabling
    aggressive kernel fusion and memory planning.
    \item \textbf{Rust-native}: No Python runtime, no GIL, no garbage
    collection pauses. Candle runs as a compiled binary with deterministic
    memory management.
    \item \textbf{Multi-backend}: CUDA (NVIDIA), Metal (Apple), Vulkan
    (cross-platform), and CPU (SIMD-optimized) backends share the same API.
    \item \textbf{Model format agnostic}: Loads ONNX, GGUF, SafeTensors, and
    native Candle format. Automatic quantization (INT8/INT4) at load time.
\end{itemize}

\subsection{Performance}

\begin{table}[H]
\centering
\small
\begin{tabular}{lrrr}
\toprule
\textbf{Model} & \textbf{Candle} & \textbf{PyTorch} & \textbf{llama.cpp} \\
\midrule
Qwen3-8B (FP16, A100) & 142 tok/s & 128 tok/s & --- \\
Qwen3-72B (INT4, A100) & 38 tok/s & 31 tok/s & 35 tok/s \\
Qwen3-8B (INT4, M4 CPU) & 28 tok/s & 11 tok/s & 26 tok/s \\
MobileNet-V3 (FP16, A100) & 0.4~ms & 0.6~ms & --- \\
Stable Diffusion (FP16, A100) & 2.8~s & 3.4~s & --- \\
\bottomrule
\end{tabular}
\caption{Inference throughput comparison. Candle's advantage comes from static
graph optimization, kernel fusion, and zero-copy tensor management.}
\label{tab:candle}
\end{table}

\subsection{Base Functions Integration}

Candle models can be invoked from Base Functions via the \texttt{candle}
module in the Goja VM. This enables on-device inference inside a Base
application without an external model serving infrastructure:

\begin{lstlisting}[language=JavaScript,caption={Candle inference from a Base Function.}]
// Base Function: classify an image
onRecordCreate((e) => {
  const img = e.record.get("image");
  const result = candle.infer("mobilenet-v3", img);
  e.record.set("label", result.top_class);
  e.record.set("confidence", result.score);
});
\end{lstlisting}

% ============================================================================
\section{Jin Architecture}
\label{sec:jin}
% ============================================================================

Jin is Hanzo's unified multimodal AI architecture~\cite{hanzojin}. Where most
AI systems treat vision, language, and audio as separate pipelines with
separate models, Jin unifies them in a single framework with shared
representations.

\subsection{Architecture}

Jin consists of three components:

\begin{enumerate}[leftmargin=1.4em]
    \item \textbf{Modality encoders}: Separate encoders for text (tokenizer +
    embedding), vision (ViT patch encoder), and audio (mel-spectrogram +
    convolutional encoder). Each encoder projects its modality into a shared
    latent space of dimension $d = 4096$.
    \item \textbf{Fusion transformer}: A standard transformer decoder that
    operates on the concatenated encoded representations. Cross-modal
    attention enables the model to attend across modalities (e.g., attending
    to an image region when generating text about it).
    \item \textbf{Modality decoders}: Separate decoders for text (token
    generation), vision (diffusion decoder for image generation), and audio
    (vocoder for speech synthesis).
\end{enumerate}

The key property: a single Jin model handles text-to-text, text-to-image,
image-to-text, audio-to-text, and any other modality combination without
separate model deployments. On Hanzo AI Cloud, a Jin model occupies one GPU
slot but serves all modality combinations.

\subsection{Mixture of Experts}

Jin uses Mixture of Experts (MoE) routing in the fusion transformer. Each
transformer block routes tokens to $k = 2$ of $E = 16$ experts, based on a
learned gating function:

\begin{equation}
\text{output} = \sum_{i \in \text{top-}k(g(\bm{x}))} g_i(\bm{x}) \cdot
\text{Expert}_i(\bm{x})
\end{equation}

MoE enables Jin to scale to 128B total parameters while activating only 16B
per token, keeping inference cost comparable to a dense 16B model. Different
experts specialize in different modality combinations, learned during training
without explicit supervision.

% ============================================================================
\section{Edge Inference}
\label{sec:edge}
% ============================================================================

Not all inference should happen in the cloud. For latency-sensitive
applications (real-time object detection, voice assistants, on-device
classification), Hanzo AI Cloud supports edge inference via Hanzo
Edge~\cite{hanzoedge}.

\subsection{On-Device Runtime}

Hanzo Edge is a Candle-based inference runtime for mobile and IoT devices. It
supports:

\begin{itemize}[leftmargin=1.1em]
    \item \textbf{Apple Neural Engine}: CoreML integration for ANE-accelerated
    inference on iPhone/iPad/Mac.
    \item \textbf{GPU compute}: Metal (Apple), Vulkan (Android/Linux) for
    GPU-accelerated inference.
    \item \textbf{CPU fallback}: NEON (ARM) and AVX-512 (x86) SIMD-optimized
    CPU inference for devices without GPU access.
    \item \textbf{Dynamic quantization}: INT8/INT4 quantization applied at
    model load time, adapting to available memory.
\end{itemize}

\subsection{CRDT Result Sync}

Edge inference results sync to the user's Base instance via the CRDT protocol
described in the Base runtime paper~\cite{hanzobaseruntime}. A mobile device
runs inference locally, writes the result to its local Base vault, and the
result syncs to the user's other devices and cloud Base instance on next
connection. This enables offline inference with eventual consistency.

\subsection{Model Distribution}

Edge models are distributed from the Hanzo model registry via a CDN-backed
download with integrity verification:

\begin{equation}
\text{verify}(\text{model}_{\text{local}}) :
\text{SHA3-256}(\text{model}_{\text{local}}) \stackrel{?}{=}
\text{model\_id}_{\text{A-Chain}}
\end{equation}

The device downloads the model once, verifies its hash against the A-Chain
registry entry, and caches it locally. Model updates are delta-compressed
(binary diff of quantized weights) to minimize bandwidth.

% ============================================================================
\section{Provider Marketplace}
\label{sec:marketplace}
% ============================================================================

Hanzo AI Cloud is not a single-operator cloud. It is a marketplace of compute
providers, each staking HANZO tokens as collateral for their service quality
commitments.

\subsection{Stake-Weighted Selection}

Providers register on-chain with:

\begin{itemize}[leftmargin=1.1em]
    \item \textbf{Hardware attestation}: Verified GPU inventory (model, VRAM,
    count) via TPM-based attestation.
    \item \textbf{HANZO stake}: Minimum 10,000 HANZO per GPU. Higher stake
    increases scheduling priority.
    \item \textbf{SLA declaration}: Committed latency (p99), uptime (monthly),
    and throughput (tokens/sec).
    \item \textbf{Pricing}: Per-token, per-image, or per-second pricing in
    HANZO.
\end{itemize}

The scheduler selects providers using a score function:

\begin{equation}
\text{score}(p) = \alpha \cdot \frac{S_p}{\sum_j S_j} +
\beta \cdot \frac{1}{L_p} + \gamma \cdot \frac{U_p}{T_p} +
\delta \cdot \frac{1}{C_p}
\end{equation}

where $S_p$ is stake, $L_p$ is historical latency, $U_p/T_p$ is uptime ratio,
$C_p$ is price, and $\alpha + \beta + \gamma + \delta = 1$ are
user-configurable weights. Users who prioritize cost set $\delta$ high; users
who prioritize reliability set $\gamma$ high.

\subsection{SLA Enforcement}

SLA violations are detected automatically and enforced on-chain:

\begin{itemize}[leftmargin=1.1em]
    \item \textbf{Latency violation}: If measured p99 exceeds declared p99 by
    $>$20\% over a 1-hour window, the provider is penalized 0.1\% of stake.
    \item \textbf{Uptime violation}: If monthly uptime falls below declared
    SLA, the provider is penalized proportionally (up to 5\% of stake per
    month).
    \item \textbf{Incorrect output}: If a verifier detects an incorrect
    inference result (via PoAI spot-checking), the provider is slashed up
    to 25\% of stake.
    \item \textbf{Compensation}: Slashed HANZO is distributed to the affected
    users as compensation, not burned.
\end{itemize}

\subsection{Compute Bidding}

For batch workloads (training, fine-tuning, bulk inference), providers submit
bids in a sealed-bid auction:

\begin{enumerate}[leftmargin=1.4em]
    \item The user submits a job specification (model, dataset size, deadline).
    \item Providers submit sealed bids (price, estimated completion time).
    \item The scheduler selects the bid with the best score (price $\times$
    time, weighted by stake and historical performance).
    \item Payment is escrowed in HANZO. Released on verified job completion.
\end{enumerate}

% ============================================================================
\section{Comparison to Existing Platforms}
\label{sec:comparison}
% ============================================================================

\begin{table}[H]
\centering
\small
\begin{tabular}{lccccc}
\toprule
\textbf{Feature} & \textbf{Hanzo} & \textbf{OpenAI} & \textbf{Anthropic} &
\textbf{Replicate} & \textbf{Together} \\
\midrule
Confidential inference & \checkmark & $\times$ & $\times$ & $\times$ & $\times$ \\
FHE support & \checkmark & $\times$ & $\times$ & $\times$ & $\times$ \\
Agent capsules & \checkmark & $\times$ & $\times$ & $\times$ & $\times$ \\
Verifiable receipts & \checkmark & $\times$ & $\times$ & $\times$ & $\times$ \\
Multi-provider & \checkmark & $\times$ & $\times$ & \checkmark & \checkmark \\
Open model registry & \checkmark & $\times$ & $\times$ & \checkmark & \checkmark \\
Edge inference & \checkmark & $\times$ & $\times$ & $\times$ & $\times$ \\
On-chain SLA & \checkmark & $\times$ & $\times$ & $\times$ & $\times$ \\
Self-hostable & \checkmark & $\times$ & $\times$ & $\times$ & $\times$ \\
User-controlled keys & \checkmark & $\times$ & $\times$ & $\times$ & $\times$ \\
\bottomrule
\end{tabular}
\caption{Platform comparison. Hanzo AI Cloud is the only platform offering
confidential inference with verifiable attestation.}
\label{tab:comparison}
\end{table}

\paragraph{OpenAI.} Market leader in model quality. No confidential inference.
No verifiable outputs. No self-hosting. Data retention policy permits 30-day
logging. Closed-source models with no weight access.

\paragraph{Anthropic.} Strong safety focus. No confidential inference. No
agent isolation beyond conversation-level context windows. No self-hosting.
Constitutional AI provides behavioral guarantees, not cryptographic ones.

\paragraph{Replicate.} Open model marketplace. No confidential inference.
Shared GPU instances with no tenant isolation. No verifiable outputs. Good
developer experience for model deployment, but no sovereignty guarantees.

\paragraph{Together.} Open model serving with competitive pricing. No
confidential inference. No agent hosting. No verifiable outputs. Together
is a cheaper alternative to OpenAI for open models, not a fundamentally
different architecture.

\paragraph{Hanzo differentiation.} Hanzo AI Cloud is the only platform where
the provider mathematically cannot see the user's data (FHE), the agent's
memory is cryptographically private (capsule vault), and every inference result
is provably attributable (A-Chain receipts). This is not a feature advantage.
It is an architectural category difference.

% ============================================================================
\section{Conclusion}
\label{sec:conclusion}
% ============================================================================

The AI cloud industry is built on a broken trust model. Users send their most
sensitive data to providers who promise---but cannot prove---that they will not
read, store, or misuse it. This trust model fails under adversarial conditions
(breaches, subpoenas, rogue employees) and is incompatible with regulated
industries that require cryptographic guarantees.

Hanzo AI Cloud replaces trust with proof:

\begin{enumerate}[leftmargin=1.4em]
    \item \textbf{Confidential inference}: CKKS FHE ensures the provider
    evaluates models on encrypted inputs. The provider never sees the
    plaintext. This is a mathematical guarantee, not a policy.
    \item \textbf{Agent sovereignty}: Each agent's vault, memory, and identity
    are cryptographically isolated. The platform cannot access agent state
    without the agent's keys.
    \item \textbf{Verifiable serving}: Every inference result carries a
    receipt anchored to the Lux A-Chain, proving provenance. Regulators,
    auditors, and users can verify exactly which model produced which output.
    \item \textbf{Decentralized provision}: No single provider controls the
    infrastructure. Stake-weighted selection and on-chain SLA enforcement
    align provider incentives with user outcomes.
    \item \textbf{Edge-to-cloud continuum}: The same model format, the same
    attestation, the same CRDT sync protocol work from a phone to a data
    center. Inference happens where it makes sense.
\end{enumerate}

The era of ``trust us, we won't look at your data'' is ending. The era of
``we cannot look at your data'' is beginning. Hanzo AI Cloud is the
infrastructure for the latter.

% ============================================================================
% REFERENCES
% ============================================================================
\begin{thebibliography}{99}

\bibitem{hanzocandle}
Z.~Kelling, D.~Wei, and M.~Chen.
\newblock Hanzo Candle: A {Rust} {ML} framework for high-performance inference.
\newblock \emph{Hanzo AI Research}, August 2024.

\bibitem{hanzofhe}
Z.~Kelling.
\newblock Hanzo {FHE} Inference: Fully homomorphic encryption for confidential
{AI} inference.
\newblock \emph{Hanzo AI Research}, April 2026.

\bibitem{hanzojin}
Z.~Kelling.
\newblock Jin: Unified multimodal {AI} architecture for vision-language, audio,
and cross-modal understanding.
\newblock \emph{Hanzo AI Research}, April 2023.

\bibitem{hanzoedge}
Z.~Kelling.
\newblock Hanzo Edge: On-device {AI} inference for mobile and {IoT} deployments.
\newblock \emph{Hanzo AI Research}, July 2024.

\bibitem{hanzobaseruntime}
Z.~Kelling.
\newblock Hanzo Base: The single-binary runtime for sovereign applications.
\newblock \emph{Hanzo AI Research}, April 2026.

\bibitem{luxwhitepaper}
{Lux Partners}.
\newblock Lux: A scalable multi-consensus blockchain with post-quantum
cryptography.
\newblock \emph{Lux Network Technical Report}, 2024.

\bibitem{hanzoplatform}
M.~Chen, D.~Wei, and Z.~Kelling.
\newblock Hanzo Platform: {PaaS} for {AI}-native application deployment.
\newblock \emph{Hanzo AI Research}, December 2021.

\bibitem{ckks2017}
J.~H.~Cheon, A.~Kim, M.~Kim, and Y.~Song.
\newblock Homomorphic encryption for arithmetic of approximate numbers.
\newblock In \emph{ASIACRYPT}, Springer, 2017.

\bibitem{openaiapi}
{OpenAI}.
\newblock {API} data usage policies.
\newblock \url{https://openai.com/policies/api-data-usage}, 2024.

\bibitem{gentry2009}
C.~Gentry.
\newblock Fully homomorphic encryption using ideal lattices.
\newblock In \emph{STOC}, ACM, 2009.

\bibitem{brakerski2014}
Z.~Brakerski, C.~Gentry, and V.~Vaikuntanathan.
\newblock ({Leveled}) Fully homomorphic encryption without bootstrapping.
\newblock \emph{ACM Transactions on Computation Theory}, 6(3):1--36, 2014.

\bibitem{gilad2016}
R.~Gilad-Bachrach, N.~Dowlin, K.~Laine, K.~Lauter, M.~Naehrig, and
J.~Wernsing.
\newblock {CryptoNets}: Applying neural networks to encrypted data with high
throughput and accuracy.
\newblock In \emph{ICML}, 2016.

\bibitem{luxagent}
Z.~Kelling.
\newblock Lux Agent Sovereignty: Decentralized identity and capability
management for autonomous agents.
\newblock \emph{Lux Network Technical Report}, 2025.

\end{thebibliography}

\end{document}
